\section{Related Work}
\label{sec:related_work}

\paragraph{Multi-task optimization via gradient and loss manipulation.}
A large body of work modifies the joint update to mitigate task conflict, either by
combining per-task gradients, PCGrad \citep{pcgrad2020}, CAGrad \citep{cagrad2021},
MGDA \citep{sener2018mgda}, and the Nash bargaining solution of Nash-MTL 
\citep{nash-mtl2022}, or by adaptively reweighting losses, as in GradNorm
\citep{chen2018gradnormgradientnormalizationadaptive}, uncertainty weighting \citep{kendall2018multitasklearningusinguncertainty}, and FAMO
\citep{liu2023famofastadaptivemultitask}, which balances progress in $O(1)$ space and time rather than storing
all per-task gradients. Simpler alternatives such as random loss weighting
\citep{lin2022reasonableeffectivenessrandomweighting} and unitary scalarization \citep{unitary2022kurin} match these
methods. \citep{unitary2022kurin} found no clear advantage of aforementioned methods against PCGrad, so we do not expect further variants to
change our conclusions.

\paragraph{Data mixture and domain reweighting.}
For pre-training, domain weights are set by proxy-model or parametric methods: DoReMi
\citep{xie2023doremioptimizingdatamixtures}, DoGE \citep{fan2024dogedomainreweightinggeneralization}, Skill-it \citep{chen2023skillitdatadrivenskillsframework},
data-mixing laws \citep{li2025datamixingoptimizationsupervised}, RegMix \citep{liu2025regmixdatamixtureregression},  and AutoScale
\citep{kang2025autoscalescaleawaredatamixing}, with online variants ODM \citep{albalak2023efficientonlinedatamixing},  ADO
\citep{jiang2024adaptivedataoptimizationdynamic}, and Aioli \citep{chen2025aioli}. Fine-tuning studies examine SFT
composition and interference \citep{dong-etal-2024-abilities}. These fit a mixing law or
bandit weights against a proxy loss. PDCM optimizes cumulative \emph{policy improvement}
online under a no-task-left-behind constraint, without a parametric mixing law.

% \subsection{Aioli}
% \label{sec:bg_constrained_allocation}

% A natural approach to task allocation is to optimize the mixture according to a model of how allocation affects subsequent training progress. Aioli and its predecessors ~\citep{chen2025aioli,ye2025datamixinglawsoptimizing,chen2023skillitdatadrivenskillsframework,xie2023doremioptimizingdatamixtures,fan2024dogedomainreweightinggeneralization} formalizes this idea through the following fomalization.
% \begin{equation}
% \label{eq:bg_lmo}
% \max_{p_1,\ldots,p_T\in\Delta^K}\ \sum_{t=1}^{T} U_t(p_t)
% \quad \text{s.t.}\quad
% U_{t,k}(p_t)
%     =
%     \phi_k \left( \sum_{j=1}^{K} A_{t,kj}p_{t,j} \right)
% \ \forall k,t.
% \end{equation}
% where $\phi_k$ specifies the shape of the scaling law. $A_t\in\mathbb{R}^{K\times K}$ models cross-task interactions. Through \Cref{eq:bg_lmo}, the current mixture is mapped to future task improvements through a parametric cross-task response matrix, and the next allocation is chosen according to these predicted gains. Aioli's allocation score for task $j$ aggregates its predicted effect across all tasks,
% $s_{t,j}=\sum_{k=1}^{K}A_{t,kj}$, leading to an exponential update of the form
% \begin{equation}
% \label{eq:bg_aioli_update}
%     p_{t+1,j}
%     =
%     \frac{
%         p_{t,j}\exp\!\left(\eta s_{t,j}\right)
%     }{
%         \sum_{\ell=1}^{K}
%         p_{t,\ell}\exp\!\left(\eta s_{t,\ell}\right)
%     }.
% \end{equation}

% This formulation makes dynamic allocation tractable by prescribing a functional relationship between allocation and future training progress. \textbf{PDCM retains the constrained-optimization view, but does not require such a parametric response law.} Instead, it uses a single-step directional concavity condition to construct first-order constraints directly from the observed task-improvement.